\begin{figure}[htbp]
    \raggedright
    \resizebox{0.98\textwidth}{!}{%
    \begin{tikzpicture}[
        font=\small,
        arrow/.style={-{Latex[length=3mm]}, thick},
        init/.style={draw, rounded corners=2mm, minimum width=4.2cm,
                    minimum height=1.22cm, align=center, thick, fill=orange!10},
        curriculum/.style={draw, rounded corners=2mm, minimum width=5.4cm,
                          minimum height=1.95cm, align=center, thick, fill=blue!7},
        final/.style={draw, rounded corners=2mm, minimum width=4.9cm,
                     minimum height=1.35cm, align=center, thick, fill=green!10},
        control/.style={draw, rounded corners=2mm, minimum width=4.9cm,
                       minimum height=1.35cm, align=center, thick, fill=gray!10},
        note/.style={font=\footnotesize, align=center}
    ]
        \node[init] (scratchinit) at (0,5.4) {Inicialización aleatoria\\arquitectura EEG-XL};
        \node[curriculum] (scratchcurr) at (6.2,5.4) {\textbf{Currículo EEG autosupervisado}\\
              lectura: Nieuwland + ZuCo 2.0\\
              $\downarrow$ mejor checkpoint\\
              escucha: SparrKULee + ds007808};
        \node[final] (scratchft) at (12.7,5.4) {Fine-tuning en ds004408\\conjunto EEG desde cero};

        \node[init] (megeeginit) at (0,2.2) {Inicialización desde MEG-XL\\pesos compatibles};
        \node[curriculum] (megeegcurr) at (6.2,2.2) {\textbf{Currículo EEG autosupervisado}\\
              embedding EEG específico\\
              lectura $\rightarrow$ escucha\\
              mismo conjunto y particiones};
        \node[final] (megeegft) at (12.7,2.2) {Fine-tuning en ds004408\\transferencia MEG-XL\\con embedding EEG};

        \node[init] (megmeginit) at (0,-1.0) {Inicialización desde MEG-XL\\pesos compatibles};
        \node[curriculum] (megmegcurr) at (6.2,-1.0) {\textbf{Currículo EEG autosupervisado}\\
              embedding MEG reutilizado\\
              lectura $\rightarrow$ escucha\\
              mismo conjunto y particiones};
        \node[final] (megmegft) at (12.7,-1.0) {Fine-tuning en ds004408\\transferencia MEG-XL\\con embedding MEG};

        \node[control] (directinit) at (0,-4.0) {Inicialización aleatoria};
        \node[control] (direct) at (6.2,-4.0) {Sin currículo autosupervisado};
        \node[control] (directft) at (12.7,-4.0) {Fine-tuning directo en ds004408\\control supervisado};

        \draw[arrow] (scratchinit) -- (scratchcurr);
        \draw[arrow] (scratchcurr) -- (scratchft);
        \draw[arrow] (megeeginit) -- (megeegcurr);
        \draw[arrow] (megeegcurr) -- (megeegft);
        \draw[arrow] (megmeginit) -- (megmegcurr);
        \draw[arrow] (megmegcurr) -- (megmegft);
        \draw[arrow] (directinit) -- (direct);
        \draw[arrow] (direct) -- (directft);

        \node[note] at (6.2,-5.45) {Las cuatro condiciones comparten el mismo conjunto de fine-tuning y el mismo protocolo de evaluación.};
    \end{tikzpicture}}
    \caption{Diseño experimental de la adaptación de MEG-XL a EEG-XL. Las tres rutas autosupervisadas mantienen el mismo currículo de lectura--escucha y difieren en la inicialización y en el embedding de tipo de sensor; el control supervisado permite separar el efecto del preentrenamiento EEG de la evaluación directa sobre ds004408.}
    \label{fig:eegxl_design}
\end{figure}
